% RACER-GT mathematical appendix (English).
\documentclass[12pt,a4paper,oneside]{article}
% English setup. Compiled with XeLaTeX so that both languages share one engine
% and one font family, which keeps the two PDFs typographically comparable.
\usepackage{fontspec}
\setmainfont{Noto Serif}
\setsansfont{Noto Sans}
\setmonofont{Noto Sans Mono}

% Author-year citations. Loaded here rather than in _common because it must
% precede hyperref, and because the Chinese manuscript uses numeric citations
% that natbib's author-year mode would reject. Bibliographies are hand-written
% thebibliography blocks: biblatex/biber is absent from a BasicTeX install, and
% a document set that only builds on a full TeX Live is not reproducible enough
% for this project.
\usepackage[authoryear,round]{natbib}

% Single source of truth for version and date across every RACER-GT document.
% Regenerate nothing by hand: bump here, rebuild, and every PDF agrees.
\newcommand{\racerversion}{1.1.0}
\newcommand{\racerdate}{2026-07-27}
\newcommand{\racerauthor}{Yi-Hao Lai}
\newcommand{\raceraffil}{Department of Finance, Da-Yeh University}
\newcommand{\raceremail}{yhlai@mail.dyu.edu.tw}
\newcommand{\racerrepo}{https://github.com/yhlai0911/RACER-GT}

% Shared package set and mathematical macros for every RACER-GT document,
% in both languages. Language-specific font and theorem-name setup lives in
% _preamble-zh.tex and _preamble-en.tex.
%
% Font policy: the build depends only on the five Noto families below. They are
% installable on Linux CI (fonts-noto-core, fonts-noto-cjk) and on macOS via
% Homebrew casks. Do not introduce a font outside this set without adding it to
% .github/workflows/docs.yml as well, or the PDF build stops being reproducible.

\usepackage[margin=2.35cm,headheight=15pt]{geometry}
\usepackage{amsmath,amssymb,amsthm,mathtools,bm}
\usepackage{booktabs,longtable,tabularx,array,multirow}
\usepackage{graphicx,float,caption,subcaption}
\usepackage{enumitem}
\usepackage{xcolor}
\usepackage{microtype}
\usepackage{setspace}
\usepackage{fancyhdr}
\usepackage{hyperref}
\usepackage[nameinlink,noabbrev]{cleveref}
\usepackage{algorithm}
\usepackage{algpseudocode}
\usepackage{listings}
\usepackage{tikz}
\usetikzlibrary{arrows.meta,positioning,shapes.geometric,fit,calc}

\hypersetup{
  colorlinks=true,
  linkcolor=black,
  citecolor=black,
  urlcolor=blue,
  pdfauthor={\racerauthor}
}

\pagestyle{fancy}
\fancyhf{}
\rhead{v\racerversion}
\cfoot{\thepage}
\setstretch{1.28}
\setlength{\parskip}{0.35em}
\setlist[itemize]{leftmargin=2em,itemsep=0.25em,topsep=0.3em}
\setlist[enumerate]{leftmargin=2.2em,itemsep=0.25em,topsep=0.3em}

% ---------------------------------------------------------------- mathematics
\newcommand{\E}{\mathbb{E}}
\newcommand{\Var}{\operatorname{Var}}
\newcommand{\Cov}{\operatorname{Cov}}
\newcommand{\tr}{\operatorname{tr}}
\newcommand{\diag}{\operatorname{diag}}
\newcommand{\argmin}{\operatorname*{arg\,min}}
\newcommand{\argmax}{\operatorname*{arg\,max}}
\newcommand{\one}{\bm{1}}
\newcommand{\Rset}{\mathbb{R}}
\newcommand{\R}{\mathbb{R}}
\newcommand{\plim}{\operatorname*{plim}}
\newcommand{\RACER}{\textsc{Racer-GT}}
\newcommand{\indic}[1]{\mathbb{I}\!\left\{#1\right\}}
\newcommand{\Sig}{\bm{\Sigma}}
\newcommand{\wvec}{\bm{w}}
\newcommand{\evec}{\bm{e}}
\newcommand{\Zvec}{\bm{Z}}

% Design facets, used identically in both languages so that a reader can move
% between the Chinese and English documents without re-learning notation.
\newcommand{\facetT}{\mathcal{T}}   % historical date
\newcommand{\facetD}{\mathcal{D}}   % collection day
\newcommand{\facetS}{\mathcal{S}}   % collection stream
\newcommand{\facetR}{\mathcal{R}}   % technical replicate

\lstset{
  basicstyle=\ttfamily\footnotesize,
  breaklines=true,
  frame=single,
  rulecolor=\color{gray!40},
  backgroundcolor=\color{gray!5},
  columns=fullflexible,
  keepspaces=true,
  showstringspaces=false
}


\setlength{\parindent}{1.5em}

\newtheorem{assumption}{Assumption}[section]
\newtheorem{proposition}{Proposition}[section]
\newtheorem{corollary}{Corollary}[section]
\newtheorem{remark}{Remark}[section]
\newtheorem{definition}{Definition}[section]
\newtheorem{lemma}{Lemma}[section]

\lhead{RACER-GT: Mathematical Appendix}

\title{\bfseries RACER-GT Mathematical Appendix\\[0.4em]
\large Derivations, Proofs, and Estimator Details}
\author{\racerauthor\\\small \raceraffil}
\date{\racerdate\ \textbullet\ version \racerversion}

\begin{document}
\maketitle

\begin{abstract}
\noindent This appendix collects the derivations that the methodology
manuscript states without full working, and documents the exact estimator each
software routine implements. It is written to be checkable: every result is
tied to the module and function that computes it, so a reader can verify that
the code and the mathematics agree. Where the implementation departs from the
textbook estimator---and it does so in three places---the departure is stated
and justified rather than left for the reader to discover.
\end{abstract}

\tableofcontents
\newpage

% =============================================================================
\section{Notation}

\begin{table}[H]
\centering
\small
\caption{Symbols used throughout.}
\begin{tabularx}{\textwidth}{llX}
\toprule
Symbol & Domain & Meaning\\
\midrule
$t$ & $1,\dots,T$ & historical calendar date\\
$i$ & $1,\dots,N$ & complete pull (one full history reconstruction)\\
$j,k$ & $1,\dots,J$ & chunk (a fixed query window)\\
$d,s,r$ & --- & collection day, stream, technical replicate\\
$Y_{jt}$ & $[0,100]$ & raw GT value, date $t$, chunk $j$\\
$a_j,\ \ell_j$ & $\R_{>0},\R$ & chunk scale and its logarithm, $\ell_j=\log a_j$\\
$X_t$ & $\R_{>0}$ & latent common-scale signal\\
$Z_{it}$ & $\R$ & calibrated value of pull $i$ at date $t$\\
$\Sig$ & $\R^{N\times N}$ & residual covariance across pulls\\
$\wvec$ & $\Delta^{N-1}$ & consensus weights, $\one^\top\wvec=1$\\
$c_{\min}$ & $\R_{>0}$ & minimum raw value admitted to an overlap edge\\
$m_{\min}$ & $\mathbb{N}$ & minimum usable overlap days for an edge\\
\bottomrule
\end{tabularx}
\end{table}

% =============================================================================
\section{Robust edge estimation}

\subsection{The estimator}

For an edge $(j,k)$, let $r_1,\dots,r_n$ be the log ratios
$\log Y_{jt}-\log Y_{kt}$ over the usable overlap days. The Huber location
estimate solves
\begin{equation}
\sum_{i=1}^{n}\psi_c\!\left(\frac{r_i-\hat{\mu}}{\hat{\sigma}}\right)=0,
\qquad
\psi_c(u)=\max\{-c,\min(c,u)\},
\end{equation}
by iteratively reweighted least squares. Writing
$w_i=\psi_c(u_i)/u_i$ for $u_i=(r_i-\mu)/\hat{\sigma}$, the update is
\begin{equation}
\mu^{(m+1)}=\frac{\sum_i w_i^{(m)} r_i}{\sum_i w_i^{(m)}},
\end{equation}
which is a weighted mean that downweights observations beyond $c\hat\sigma$.
The scale $\hat{\sigma}$ is the normalized median absolute deviation.

\subsection{Variance of the edge estimate}

The asymptotic variance of an M-estimator of location is
\begin{equation}
\Var(\hat{\mu}) \;\approx\;
\frac{\hat{\sigma}^2}{n}\cdot
\frac{\E[\psi_c(u)^2]}{\big(\E[\psi_c'(u)]\big)^2}.
\label{eq:m-var}
\end{equation}
The implementation uses the finite-sample weighted counterpart: with the final
Huber weights $w_i=\psi_c(u_i)/u_i$ it forms an effective sample size and divides
the robust scale by it,
\begin{equation}
n_{\text{eff}}=\frac{\left(\sum_i w_i\right)^2}{\sum_i w_i^2},
\qquad
\hat{v}_{jk}=\frac{\hat{\sigma}^2_{\text{robust}}}{\max(n_{\text{eff}},1)}.
\end{equation}
The two agree closely under normality at $c=1.345$, where the efficiency factor
is about 1.04, but they are not the same estimator: \eqref{eq:m-var} is
asymptotic, whereas the weighted form responds directly to a handful of
observations dominating a single edge. \texttt{huber\_location} in
\texttt{overlap.py} computes the latter.
This variance becomes the edge precision $1/\hat{v}_{jk}$ in the graph problem.
Two implementation details matter:
\begin{enumerate}
\item a floor is applied, $\hat{v}\leftarrow\max(\hat{v},v_{\text{floor}})$,
because an edge with near-zero estimated variance would otherwise dominate the
normal equations entirely;
\item a ceiling is applied to the resulting weight, $1/\hat{v}\le w_{\max}$,
for the same reason from the other direction.
\end{enumerate}

\begin{remark}[Departure 1 from the textbook estimator]
The floor and ceiling are not part of the classical M-estimation theory. They
are numerical guards. Their effect is to bound the influence any single edge can
exert on the global solution, which matters because an edge computed from
exactly $m_{\min}$ nearly identical days can produce a spuriously small
variance.
\end{remark}

% =============================================================================
\section{Graph weighted least squares}

\subsection{Normal equations}

With $B$ the reduced incidence matrix (reference column deleted),
$\bm{r}$ the stacked edge estimates and $\Omega=\diag(\hat{v}_e)$,
\begin{equation}
(B^\top\Omega^{-1}B)\,\hat{\bm{\ell}}=B^\top\Omega^{-1}\bm{r}.
\end{equation}
The matrix $L=B^\top\Omega^{-1}B$ is the weighted graph Laplacian with the
reference row and column removed, sometimes called the grounded Laplacian.

\begin{proposition}[Positive definiteness]
If $G$ is connected then the grounded Laplacian $L$ is positive definite.
\end{proposition}

\begin{proof}
For any $\bm{x}\neq\bm{0}$ on the free nodes, extend it to $\tilde{\bm{x}}$ by
setting the reference entry to zero. Then
$\bm{x}^\top L\bm{x}=\sum_{e=(j,k)}\hat{v}_e^{-1}(\tilde{x}_j-\tilde{x}_k)^2\ge 0$,
with equality only if $\tilde{\bm{x}}$ is constant on every connected component.
Connectivity forces one component, and the reference entry pins that constant to
zero, so $\tilde{\bm{x}}=\bm{0}$, contradicting $\bm{x}\neq\bm{0}$.
\end{proof}

\subsection{Diagnostics that follow from the residual}

The fitted residual $\bm{\varepsilon}=\bm{r}-B\hat{\bm{\ell}}$ is exactly the
cycle inconsistency of the edge measurements: for any cycle in $G$, the sum of
edge estimates around it should be zero, and $\bm{\varepsilon}$ measures the
failure. The implementation reports:
\begin{itemize}
\item connectivity and the number of components;
\item the condition number of $L$, which flags a graph held together by a
single weak edge;
\item the weighted residual sum of squares as a cycle-consistency statistic;
\item per-node standard errors from $\diag(L^{-1})$.
\end{itemize}
None of these quantities exists under sequential stitching, because sequential
stitching uses a spanning path and therefore has no cycles to be inconsistent
about.

\subsection{Lognormal bias correction}

What is corrected is the Jensen bias \emph{of the estimator}, not the
median-versus-mean gap of the multiplicative error in the data-generating
process. With $\hat{\ell}_j$ approximately normal and
$s_j^2=\Var(\hat{\ell}_j)$,
\begin{equation}
\E\!\left[\exp(-\hat{\ell}_j)\right]
=\exp(-\ell_j)\exp\!\left(\tfrac{1}{2}s_j^2\right),
\end{equation}
so back-transforming with $\exp(-\hat{\ell}_j)$ alone systematically overstates
$a_j^{-1}$. The first-order correction is
\begin{equation}
\widehat{a_j^{-1}}_{\,bc}=\exp\!\left(-\hat{\ell}_j-\tfrac{1}{2}s_j^2\right),
\end{equation}
as in the manuscript. The implementation takes $s_j^2$ from the diagonal of
$L^{-1}$ in the graph WLS solution, which is \texttt{log\_scale\_variance} in
\texttt{overlap.py}.

\begin{remark}[Departure 2]
This correction is exact only when $\hat{\ell}_j$ is approximately normal, and
its magnitude shrinks toward zero as the number of usable overlap days grows.
That distinguishes it from the variance of the observation error
$\epsilon_{jt}$ itself, which does not shrink with sample size; the two are easy
to conflate, so the quantity actually used is named explicitly here. It is
enabled by default because the multiplicative model motivates it, and it remains
a configuration flag because normality is not testable at these sample sizes.
\end{remark}

% =============================================================================
\section{Variance components}

\subsection{Expected mean squares}

For the fully crossed design of the manuscript with $n_t,n_d,n_s$ levels and
$n_r$ replicates, the method-of-moments equations follow from the expected mean
squares. Writing $\sigma^2_{\bullet}$ for the component of the indicated effect,
\begin{align}
\E[\mathrm{MS}_T] &= \sigma^2_E + n_r\sigma^2_{TDS} + n_rn_s\sigma^2_{TD}
+ n_rn_d\sigma^2_{TS} + n_rn_dn_s\sigma^2_T,\\
\E[\mathrm{MS}_{TD}] &= \sigma^2_E + n_r\sigma^2_{TDS} + n_rn_s\sigma^2_{TD},\\
\E[\mathrm{MS}_{TDS}] &= \sigma^2_E + n_r\sigma^2_{TDS},\\
\E[\mathrm{MS}_E] &= \sigma^2_E,
\end{align}
with the analogous equations for $D$, $S$, $TS$ and $DS$. Solving downward from
the highest-order term gives each component.

\begin{remark}[Departure 3: negative components]
Method-of-moments estimates can be negative, which is impossible for a variance.
The implementation clips them at zero by default. Clipping is the standard
practical choice but it introduces a small upward bias in the clipped component
and a corresponding distortion in any ratio that uses it. The raw, unclipped
values are retained in the output so that the reader can see how often clipping
was necessary; frequent clipping is evidence that the design is too small for
the decomposition being attempted.
\end{remark}

\subsection{Generalizability and dependability coefficients}

For relative decisions (rank ordering across dates),
\begin{equation}
E\rho^2=\frac{\sigma^2_T}
{\sigma^2_T+\dfrac{\sigma^2_{TD}}{n_d}+\dfrac{\sigma^2_{TS}}{n_s}
+\dfrac{\sigma^2_{TDS}}{n_dn_s}+\dfrac{\sigma^2_E}{n_dn_s n_r}} .
\label{eq:erho-en}
\end{equation}
For absolute decisions the dependability coefficient $\Phi$ additionally charges
the main effects of $D$ and $S$, because a uniform shift matters when the score
is interpreted on its own rather than relative to other dates:
\begin{equation}
\Phi=\frac{\sigma^2_T}
{\sigma^2_T+\dfrac{\sigma^2_D}{n_d}+\dfrac{\sigma^2_S}{n_s}
+\dfrac{\sigma^2_{TD}}{n_d}+\dfrac{\sigma^2_{TS}}{n_s}
+\dfrac{\sigma^2_{DS}}{n_dn_s}
+\dfrac{\sigma^2_{TDS}}{n_dn_s}+\dfrac{\sigma^2_E}{n_dn_sn_r}} .
\label{eq:phi-en}
\end{equation}

\begin{remark}[Why $\sigma^2_{TDS}$ and $\sigma^2_E$ carry different divisors]
The $TDS$ interaction is averaged only over day $\times$ stream cells, so it is
divided by $n_dn_s$; pure replication error is averaged $n_r$ further times
within each cell, so it is divided by $n_dn_sn_r$. Collapsing the two into a
single $(\sigma^2_{TDS}+\sigma^2_E)/(n_dn_sn_r)$ term is correct only at
$n_r=1$; for $n_r>1$ it understates relative error variance and therefore
overstates $G$. The confirmatory design recommended in the manuscript uses
$n_r=2$, which is exactly the case where the difference bites.
\cref{eq:erho-en} matches \texttt{generalizability\_coefficients} in
\texttt{gstudy.py} term by term, and a consistency test in
\texttt{tests/test\_gstudy.py} keeps the two from drifting apart.
\end{remark}
The D-study evaluates both on a grid of $(n_d,n_s,n_r)$, which is what turns the
decomposition into an actionable design recommendation: the partial derivative
with respect to $n_d$ versus $n_s$ says where the next unit of effort buys more
reliability.

% =============================================================================
\section{Consensus weights under constraints}

The unconstrained solution $\wvec^*=\Sig^{-1}\one/(\one^\top\Sig^{-1}\one)$ is
derived in the manuscript. The default configuration instead solves
\begin{equation}
\min_{\wvec}\ \wvec^\top\hat{\Sig}\wvec
\quad\text{s.t.}\quad
\one^\top\wvec=1,\quad
0\le w_i\le\bar{w}.
\label{eq:qp}
\end{equation}

\begin{proposition}[The constrained problem is well posed]
If $\hat{\Sig}\succeq 0$ and $\bar{w}\ge 1/N$, the feasible set of
\eqref{eq:qp} is non-empty, compact and convex, so a minimizer exists. If
$\hat{\Sig}\succ 0$ it is unique.
\end{proposition}

\begin{proof}
The equal-weight vector $w_i=1/N$ is feasible when $\bar{w}\ge 1/N$, so the set
is non-empty; it is the intersection of a hyperplane with a box, hence compact
and convex. A continuous function on a non-empty compact set attains its
minimum. Strict convexity of the objective under $\hat{\Sig}\succ 0$ gives
uniqueness.
\end{proof}

\begin{remark}[Why $\bar{w}$ exists at all]
Without a cap, an estimated covariance that happens to make one pull look almost
noiseless will concentrate nearly all weight on it, and the consensus degenerates
toward a single pull---the very failure mode the framework exists to avoid. The
cap trades a small amount of theoretical efficiency for the guarantee that the
consensus always aggregates.
\end{remark}

\subsection{Shrinkage covariance}

The Ledoit--Wolf estimator is
$\hat{\Sig}=(1-\alpha)S+\alpha F$, with $S$ the sample covariance, $F$ a
well-conditioned target, and $\alpha\in[0,1]$ chosen to minimize expected
squared Frobenius loss. It is used because $N$ is typically small relative to
the number of dates and $S^{-1}$ is correspondingly unstable.

\subsection{What the residual covariance can resolve}\label{app:resolution}

$\hat{\Sig}$ is estimated from residuals taken about the daily central tendency of
the pulls themselves. The subtrahend is therefore correlated with what it is
subtracted from, and that bounds how different the estimated variances can be,
before any data is seen.

\begin{proposition}[Resolution bound of the residual covariance]
\label{prop:resolution}
Let $Z_{tj}=X_t+e_{tj}$ with $e_{tj}$ independent across $j$ and
$\Var(e_{tj})=\sigma_j^2$, and let residuals be taken about the cross-pull mean,
$r_{tj}=Z_{tj}-N^{-1}\sum_k Z_{tk}$. Then, writing $S=\sum_k\sigma_k^2$,
\begin{equation}
\Var(r_{tj})=\sigma_j^2\left(1-\frac{2}{N}\right)+\frac{S}{N^2}.
\label{eq:residvar}
\end{equation}
If $\sigma_2=\dots=\sigma_N=\sigma$ and $\sigma_1\to 0$, then
\begin{equation}
\frac{\Var(r_{t2})}{\Var(r_{t1})}\longrightarrow
\frac{N^2-N-1}{N-1}=N-\frac{1}{N-1},
\label{eq:resolutionbound}
\end{equation}
a limit that does not depend on $\sigma$. At $N=2$ the ratio equals one
identically, for every $\sigma_1$ and $\sigma_2$.
\end{proposition}

\begin{proof}
Write $\bar{e}_t=N^{-1}\sum_k e_{tk}$, so $r_{tj}=e_{tj}-\bar{e}_t$. By
independence $\Cov(e_{tj},\bar{e}_t)=\sigma_j^2/N$ and $\Var(\bar{e}_t)=S/N^2$,
hence
$\Var(r_{tj})=\sigma_j^2-2\sigma_j^2/N+S/N^2$, which is \eqref{eq:residvar}.
Under the stated configuration $S\to(N-1)\sigma^2$, so
$\Var(r_{t1})\to(N-1)\sigma^2/N^2$ and
$\Var(r_{t2})\to\sigma^2\bigl[(N-2)/N+(N-1)/N^2\bigr]
=\sigma^2(N^2-N-1)/N^2$. The ratio of the two is \eqref{eq:resolutionbound}, and
$\sigma^2$ cancels. For $N=2$ the factor $1-2/N$ vanishes, so \eqref{eq:residvar}
gives $\Var(r_{t1})=\Var(r_{t2})=S/4$ regardless of $\sigma_1,\sigma_2$; more
strongly, $r_{t1}=(e_{t1}-e_{t2})/2=-r_{t2}$ pointwise.
\end{proof}

\begin{remark}[What the bound is about]
\Cref{prop:resolution} concerns centring on the sample, not Google Trends. It
applies to any covariance estimated from residuals about a statistic of the same
observations, and it caps the ratio at roughly the number of observations being
combined however large the true ratio is. Simulation at a true variance ratio of
$4\times 10^{8}$ reproduces \eqref{eq:resolutionbound} to within $0.6\%$ for
$N=3,\dots,10$.
\end{remark}

\begin{remark}[The median subtrahend]
The implementation subtracts the daily median rather than the mean, which is not
a linear functional and admits no corresponding closed form. Numerically it is
better by 3 to 50 percent over the same range of $N$ and converges to the same
bound as $N$ grows, so it changes the constant and not the conclusion.
\end{remark}

\subsection{Effective number of pulls}

Two summaries are reported:
\begin{equation}
N_{\text{Kish}}=\frac{1}{\sum_i w_i^2},
\qquad
N_{\text{spec}}=\frac{\left(\sum_{m}\lambda_m\right)^2}{\sum_{m}\lambda_m^2},
\end{equation}
where $\lambda_m$ are the eigenvalues of the residual correlation matrix. The
first measures weight concentration; the second measures how much of the
variation is shared. Both equal $N$ only when weights are equal and residuals
are uncorrelated, and both fall well below $N$ in the presence of
near-duplicates. Reporting $N$ itself as the sample size would overstate
precision, which is why the acceptance rule is stated in terms of
$N_{\text{spec}}$.

\subsection{Per-day standard errors}\label{app:dailyse}

The consensus series is published with a per-day standard error. It is not a
per-day variance estimate, and because the downstream corrections of
\cref{app:eiv} consume it as though it were, the construction is given in full.

Let $\hat{X}_t$ denote the consensus on day $t$ and $y_{tj}$ the aligned value of
pull $j$. With $\hat{\Sig}$ and $\wvec$ as above, define
\begin{equation}
s_{\text{base}}=\sqrt{\wvec^{\top}\hat{\Sig}\wvec},
\qquad
M_t=\operatorname{med}_j\bigl|y_{tj}-\hat{X}_t\bigr|,
\qquad
\bar{M}=\operatorname{med}_t M_t .
\end{equation}
The reported quantity is
\begin{equation}
\mathrm{se}_t
= s_{\text{base}}\cdot
\operatorname{clip}\!\left(\frac{M_t}{\bar{M}},\;c_{\text{lo}},\;c_{\text{hi}}\right),
\qquad (c_{\text{lo}},c_{\text{hi}})=(0.25,\,4),
\label{eq:dailyse}
\end{equation}
reducing to $\mathrm{se}_t=s_{\text{base}}$ when the local term is disabled. The
published interval is $\hat{X}_t\mp 1.96\,\mathrm{se}_t$.

\begin{remark}[Three understatements, all in the same direction]
\Cref{eq:dailyse} is a conditional \emph{scale} on the dispersion across pulls,
not an estimate of $\Var(\hat{X}_t)$ day by day. Three separate features push it
downward, and none of them is a rounding matter:

\begin{enumerate}
\item \emph{The upper clip.} A day whose dispersion is genuinely ten times the
typical day is reported at four times it. The bound binds exactly on the days
where the uncertainty is largest.

\item \emph{Local reweighting is not carried into the variance.} When some pulls
are unavailable on day $t$, the consensus is formed with weights renormalized on
the available set, but $s_{\text{base}}$ is computed once from the global
$\wvec$. Concentrating weight on fewer pulls raises the true variance, and
\eqref{eq:dailyse} does not respond unless $M_t$ happens to rise with it.

\item \emph{Residuals are taken about the daily median of the pulls
themselves.} The same operation that biases $N_{\text{spec}}$ downward also
shrinks the residual dispersion from which $\hat{\Sig}$ is estimated, so
$s_{\text{base}}$ inherits that shrinkage.
\end{enumerate}

The direction deserves emphasis by contrast with $N_{\text{spec}}$, whose bias is
conservative: it understates information and therefore errs toward rejecting a
batch. These three err the other way. They make the series look more precise than
the design can support, which is the direction that flatters any conclusion drawn
from it.
\end{remark}

\begin{remark}[Consequence for the error-corrected estimators]
The moment correction of \cref{app:moment} and the error-corrected factor model
both treat $\mathrm{se}_t^2$ as a known $\Omega$. An $\Omega$ that is uniformly
too small under-corrects while every published diagnostic remains internally
consistent, so the output still reads as corrected. Detecting that state requires
the true $\Omega$, which is precisely what is unavailable. Until
\eqref{eq:dailyse} is replaced by a genuine per-day variance, results that depend
on the magnitude of $\Omega$ --- as opposed to its existence --- should be read as
conditional on a lower bound.
\end{remark}

% =============================================================================
\section{Temporal benchmarking}

\subsection{Solution}

With $Q\succ 0$, $W\succeq 0$, and $A$ the day-to-period aggregation matrix,
\begin{equation}
\bm{x}^*=(Q+A^\top WA)^{-1}(Q\bm{z}+A^\top W\bm{b}).
\end{equation}
The matrix $Q$ is built as $(\lambda_I+\lambda_R)I+\lambda_D D_1^\top D_1$ where
$D_1$ is the \emph{first}-difference operator: \texttt{\_difference\_penalty} in
\texttt{benchmark.py} forms the $(n-1)\times n$ difference matrix with offsets
$[0,1]$ and returns $D_1^\top D_1$. The first term penalizes departure from the
daily consensus, with $\lambda_R$ the ridge that keeps $Q$ positive definite when
$\lambda_I=0$; the second penalizes roughness in the correction path. Movement
preservation is thus imposed on the correction rather than on the level. A
first-difference penalty charges the \emph{slope} of the correction path; a
second-difference penalty would charge its curvature, which is a different
family of benchmarking estimators and is not what this implementation solves.

\subsection{Exact mode}

In exact mode the aggregate constraint $A\bm{x}=\bm{b}$ is imposed rather than
penalized, giving the KKT system
\begin{equation}
\begin{bmatrix} Q & A^\top\\ A & 0\end{bmatrix}
\begin{bmatrix}\bm{x}\\ \bm{\lambda}\end{bmatrix}
=
\begin{bmatrix}Q\bm{z}\\ \bm{b}\end{bmatrix},
\end{equation}
which is solvable when $A$ has full row rank. Soft mode is the default because
the lower-frequency GT series is itself measured with error, and imposing it
exactly would transfer that error into the daily series without acknowledgement.

% =============================================================================
\section{Measurement-error correction}\label{app:eiv}

\subsection{Moment correction}\label{app:moment}

Let $W_t=X_t+u_t$, let $M$ be the residual maker for the controls, and let
$\Omega_u$ be the measurement-error covariance. From
$\E[\bm{W}^\top M\bm{W}]=\E[\bm{X}^\top M\bm{X}]+\tr(M\Omega_u)$ and
$\E[\bm{W}^\top M\bm{Y}]=\beta\,\E[\bm{X}^\top M\bm{X}]$,
\begin{equation}
\hat{\beta}=\frac{\bm{W}^\top M\bm{Y}}
{\bm{W}^\top M\bm{W}-\tr(M\Omega_u)}.
\end{equation}

\begin{remark}[The denominator is a diagnostic]
A non-positive denominator means the estimated measurement error accounts for
all of the observed regressor variation after partialling out the controls. The
implementation raises an error instead of returning a number. Returning a
sign-flipped or explosive coefficient in that situation would be worse than
failing, because the failure is informative: it says the constructed index has
no usable signal left for this regression.
\end{remark}

\subsection{SIMEX}

For a grid $\lambda\in\{0,\lambda_1,\dots,\lambda_K\}$, generate
$W_t(\lambda,b)=W_t+\sqrt{\lambda}\,\tilde{u}_{tb}$ with
$\tilde{u}_{tb}$ drawn to match $\Omega_u$, estimate $\hat\beta(\lambda,b)$ for
$b=1,\dots,B$, and average over $b$. Fit a parametric curve
$g(\lambda;\bm{\theta})$ to $\{(\lambda,\bar{\beta}(\lambda))\}$ and report
$g(-1;\hat{\bm{\theta}})$, which extrapolates to the hypothetical case of no
measurement error. SIMEX is reported alongside the moment correction rather than
instead of it: agreement between the two is evidence that the classical error
structure is adequate, and disagreement is evidence that it is not.

% =============================================================================
\section{Correspondence between mathematics and code}

\begin{table}[H]
\centering
\small
\caption{Where each result is implemented.}
\begin{tabularx}{\textwidth}{llX}
\toprule
Result & Module & Entry point\\
\midrule
Huber edge estimate & \texttt{overlap.py} & \texttt{huber\_location}\\
Graph WLS, diagnostics & \texttt{overlap.py} & \texttt{OverlapGraphCalibrator.fit}\\
Exact/near duplicates & \texttt{duplicates.py} & \texttt{diagnose\_duplicates}\\
Variance components, D-study & \texttt{gstudy.py} & \texttt{run\_gstudy}\\
Design-based estimator & \texttt{consensus.py} & \texttt{fit\_design\_weighted\_consensus}\\
Constrained GLS consensus & \texttt{consensus.py} & \texttt{fit\_gls\_consensus}\\
Benchmark solution & \texttt{benchmark.py} & \texttt{temporal\_benchmark}\\
Reliability, convergence & \texttt{reliability.py} & \texttt{assess\_reliability}\\
Acceptance decision tree & \texttt{decision.py} & \texttt{evaluate\_batch}\\
Moment correction, SIMEX & \texttt{eiv.py} & \texttt{reliability\_corrected\_ols}, \texttt{simex\_ols}\\
\bottomrule
\end{tabularx}
\end{table}

\end{document}
